\chapter{Conclusion}
\label{sec:conclusion}

\section{Summary of Contributions}
\label{sec:conclusion_summary}

In this dissertation, I presented a context-aware, real-time privacy framework for
egocentric video that selectively anonymizes bystanders while preserving the faces
of active interaction partners. The framework operates as a five-stage sequential
pipeline: face detection via YOLOv8n-face, multi-face tracking via SORT, interaction
cue extraction, heuristic classification, and selective Gaussian blur anonymization.
All stages execute on CPU-only hardware without GPU acceleration, biometric face
recognition, or cloud offloading.

The primary contribution of my research is the \emph{multi-cue interaction
classifier}, which fuses four complementary signals derived from two information
sources. The first two cues (face area ratio and center distance) are extracted
from bounding-box geometry alone and require no additional model inference. The
second two cues (head-pose yaw via PnP and Mouth Aspect Ratio via MediaPipe Face
Landmarker) require facial landmark extraction but provide qualitatively different
information: yaw provides a hard veto for faces that are turned away from the
camera, while MAR-based V-VAD provides a temporal boost for faces exhibiting
sustained lip movement. The classifier combines these cues through a priority-ordered
decision function with temporal hysteresis and interaction latching, achieving PSR
= 87.50\% and IPR = 84.62\% in the Phase~2a geometry-only configuration, evaluated
against a ground-truth dataset of 66 unique face tracks from 21 egocentric video clips.
With landmark cues enabled (Phase~2b), a blind four-annotator evaluation of 29 faces
gave PSR = 100\% and IPR = 57.14\%, and a larger blind evaluation of 102 video
clips by five annotators with almost perfect agreement ($\kappa = 0.887$) confirmed this
with PSR = 96.97\% and IPR = 56.52\%. For recorded video, an offline second pass
that applies the latch retroactively and recognises the members of the wearer's
group raised IPR to 78.26\% at PSR = 93.94\% (F$_1$ = 86.40\%).

A secondary contribution is the \emph{systematic comparison with the blanket-blurring
baseline}. The baseline achieves PSR = 100\% but IPR = 0\%, making it unsuitable for
any use case in which the wearer needs to interact with people in the scene. The
proposed system recovers 84.62\% of interaction-partner utility (IPR) at a cost of
12.5 percentage points of PSR relative to the baseline. This trade-off is
appropriate for a system whose primary objective is to protect bystanders without
degrading the wearer's social experience. With the offline second pass, the
system preserved 78.26\% of the people in the wearer's group while still blurring
93.94\% of bystanders, demonstrating that selective anonymization can recover most
of the social value that blanket blurring destroys while keeping strong bystander
protection.

A third contribution is the \emph{edge-deployable implementation}. The pipeline is
designed from the outset to run on commodity CPU hardware. The measured throughput
of 11.62~FPS (Phase~2a) and approximately 9.4~FPS (Phase~2b) on a laptop CPU, with
estimated Raspberry Pi~5 throughput of 3.3 and 2.7~FPS respectively,
demonstrates that the approach is feasible on CPU-only edge hardware, and the
per-stage analysis identifies a clear path to higher frame rates through model
quantisation and hardware acceleration.

Finally, the \emph{evaluation methodology} that I developed for this project,
which includes blind annotation of still frames and of short video clips with
sound, inter-annotator agreement measurement and verification of the system's
decisions against its rendered output, forms a reusable protocol for future
egocentric privacy research.

\section{Relation to Prior Work}
\label{sec:conclusion_prior}

Compared to the most relevant prior system, BystandAR \citep{bystandar} (PSR~98.14\%,
IPR~96.27\%), the proposed framework achieves lower accuracy but eliminates the
hardware dependency on eye gaze trackers and directional microphones. Compared to
blanket-blurring systems such as EgoBlur \citep{egoblur2023} and SafeAR
\citep{ref:safear2025}, the proposed framework introduces the context-aware
classification layer that all prior blanket systems lack.

To my knowledge, the proposed system is the first to combine:
(i) real-time selective anonymization in egocentric video,
(ii) CPU-only edge deployment without biometric recognition,
(iii) a multi-cue fusion of geometry, head pose, and V-VAD, and
(iv) a quantitative evaluation distinguishing PSR from IPR as separate privacy and
utility metrics.

\section{Future Work}
\label{sec:conclusion_future}

I identify several directions for future investigation.

\textbf{Learned classifier.} The heuristic decision function could be replaced with
a lightweight learned binary classifier (e.g., a shallow MLP or logistic regression
model trained on the cue vectors) that generalises better across environments and
interaction styles than hand-tuned thresholds.

\textbf{Audio-visual fusion.} Incorporating a microphone-based voice activity
detector alongside the visual V-VAD cue would resolve ambiguous lip-movement cases
and substantially improve IPR in noisy environments.

\textbf{Independent and larger-scale evaluation.} The video-clip protocol should
be applied to a fresh set of clips, to test the offline second pass independently,
and to a larger, more diverse dataset, so that all configurations can be compared on
a common ground truth.

\textbf{Group recognition in live operation.} The retroactive latch and the
friend-group rule currently need the whole recording. Delaying the output stream
by a few seconds would give the system time to confirm partners and group members
before their frames are shown, bringing most of the benefit of the second pass to
live operation.

\textbf{Adaptive proximity threshold.} Most Phase~2b false negatives fell below the
fixed area threshold. Relaxing $\tau_{\mathrm{area}}$ when landmark cues confirm a
frontal, speaking face, and improving landmark coverage on small faces, are
direct routes to higher IPR without weakening PSR.

\textbf{On-device deployment.} Direct benchmarking on the Raspberry Pi~5 with ONNX
INT8 quantised models, combined with the input-resolution and batching optimisations
described in Section~\ref{sec:disc_fps}, would validate or refine the throughput projection and
identify any architecture-specific performance bottlenecks not captured by the
current scaling-factor estimate.

\textbf{Consent and transparency mechanisms.} The current system operates
invisibly from the bystander's perspective. Future work could integrate BystandARIA's
visible-light communication approach \citep{ref:bystandaria} as a complementary
consent channel, allowing bystanders to actively signal their privacy preferences
to the system.
